\section{Implementation}
\subsection{Hardware}
For the hardware implementation of TinyNet, we developed compact routers powered by the ATSAMS70 microcontroller from Atmel/Microchip, which operates at a core clock speed of 300~MHz. The router schematic and board file are illustrated in Figs. \ref{fig:router-schematic} and \ref{fig:router-board-file} in the appendix. While this microcontroller is capable of supporting UART transmission speeds up to 18.75 Mbps, we observed that a reduction in bitrate to 3.125 Mbps was necessary to achieve stable and reliable data transmission. At higher bitrates, the inter-byte arrival interval becomes shorter than the UART interrupt service routine latency on the ATSAMS70, causing receive buffer overruns; the 3.125 Mbps setting provides sufficient margin for the 300 MHz core to service each incoming byte interrupt before the next byte arrives.

To enhance the system's resilience to electromagnetic interference (EMI), we incorporated the ISL3177 differential driver into the UART lines. Although this measure was not strictly required for our initial testing environment, it is a prudent preparatory step toward operating TinyNet in more challenging EMI conditions, such as those found near Brushless Motor Controllers. This consideration forms a part of our agenda for future development and testing, ensuring that TinyNet remains robust and reliable in a wide range of operational settings.

\subsection{Simulation}
We implemented TinyNet's software simulation in JavaScript, running on a commodity Windows laptop. The underlying framework is adapted from Simbit, a discrete-event P2P network simulator originally designed for blockchain research~\cite{simbit}; we replaced its packet format and propagation model with the UART timing equations (equations~(\ref{eq:d})--(\ref{eq:dtx})) and defined topology as an array of adjacency lists. Each element represents the port connections of one node, allowing us to model arbitrary topologies (evaluated here on up to 16 nodes).

In the simulation, each node is represented as a ``manager.'' These managers handle incoming packets on each port, maintain their own lookup and buffer depth tables, and apply the forwarding algorithm to each incoming packet. Events are scheduled by a discrete-event simulation engine that advances the virtual clock between scheduled callbacks. We can also inject custom traffic patterns to explore TinyNet's behavior under specific conditions.